% Options for packages loaded elsewhere
\PassOptionsToPackage{unicode}{hyperref}
\PassOptionsToPackage{hyphens}{url}
\documentclass[
  12pt,
]{article}
\usepackage{xcolor}
\usepackage[margin=2.5cm]{geometry}
\usepackage{amsmath,amssymb}
\setcounter{secnumdepth}{5}
\usepackage{iftex}
\ifPDFTeX
  \usepackage[T1]{fontenc}
  \usepackage[utf8]{inputenc}
  \usepackage{textcomp} % provide euro and other symbols
\else % if luatex or xetex
  \usepackage{unicode-math} % this also loads fontspec
  \defaultfontfeatures{Scale=MatchLowercase}
  \defaultfontfeatures[\rmfamily]{Ligatures=TeX,Scale=1}
\fi
\usepackage{lmodern}
\ifPDFTeX\else
  % xetex/luatex font selection
\fi
% Use upquote if available, for straight quotes in verbatim environments
\IfFileExists{upquote.sty}{\usepackage{upquote}}{}
\IfFileExists{microtype.sty}{% use microtype if available
  \usepackage[]{microtype}
  \UseMicrotypeSet[protrusion]{basicmath} % disable protrusion for tt fonts
}{}
\makeatletter
\@ifundefined{KOMAClassName}{% if non-KOMA class
  \IfFileExists{parskip.sty}{%
    \usepackage{parskip}
  }{% else
    \setlength{\parindent}{0pt}
    \setlength{\parskip}{6pt plus 2pt minus 1pt}}
}{% if KOMA class
  \KOMAoptions{parskip=half}}
\makeatother
\usepackage{color}
\usepackage{fancyvrb}
\newcommand{\VerbBar}{|}
\newcommand{\VERB}{\Verb[commandchars=\\\{\}]}
\DefineVerbatimEnvironment{Highlighting}{Verbatim}{commandchars=\\\{\}}
% Add ',fontsize=\small' for more characters per line
\usepackage{framed}
\definecolor{shadecolor}{RGB}{248,248,248}
\newenvironment{Shaded}{\begin{snugshade}}{\end{snugshade}}
\newcommand{\AlertTok}[1]{\textcolor[rgb]{0.94,0.16,0.16}{#1}}
\newcommand{\AnnotationTok}[1]{\textcolor[rgb]{0.56,0.35,0.01}{\textbf{\textit{#1}}}}
\newcommand{\AttributeTok}[1]{\textcolor[rgb]{0.13,0.29,0.53}{#1}}
\newcommand{\BaseNTok}[1]{\textcolor[rgb]{0.00,0.00,0.81}{#1}}
\newcommand{\BuiltInTok}[1]{#1}
\newcommand{\CharTok}[1]{\textcolor[rgb]{0.31,0.60,0.02}{#1}}
\newcommand{\CommentTok}[1]{\textcolor[rgb]{0.56,0.35,0.01}{\textit{#1}}}
\newcommand{\CommentVarTok}[1]{\textcolor[rgb]{0.56,0.35,0.01}{\textbf{\textit{#1}}}}
\newcommand{\ConstantTok}[1]{\textcolor[rgb]{0.56,0.35,0.01}{#1}}
\newcommand{\ControlFlowTok}[1]{\textcolor[rgb]{0.13,0.29,0.53}{\textbf{#1}}}
\newcommand{\DataTypeTok}[1]{\textcolor[rgb]{0.13,0.29,0.53}{#1}}
\newcommand{\DecValTok}[1]{\textcolor[rgb]{0.00,0.00,0.81}{#1}}
\newcommand{\DocumentationTok}[1]{\textcolor[rgb]{0.56,0.35,0.01}{\textbf{\textit{#1}}}}
\newcommand{\ErrorTok}[1]{\textcolor[rgb]{0.64,0.00,0.00}{\textbf{#1}}}
\newcommand{\ExtensionTok}[1]{#1}
\newcommand{\FloatTok}[1]{\textcolor[rgb]{0.00,0.00,0.81}{#1}}
\newcommand{\FunctionTok}[1]{\textcolor[rgb]{0.13,0.29,0.53}{\textbf{#1}}}
\newcommand{\ImportTok}[1]{#1}
\newcommand{\InformationTok}[1]{\textcolor[rgb]{0.56,0.35,0.01}{\textbf{\textit{#1}}}}
\newcommand{\KeywordTok}[1]{\textcolor[rgb]{0.13,0.29,0.53}{\textbf{#1}}}
\newcommand{\NormalTok}[1]{#1}
\newcommand{\OperatorTok}[1]{\textcolor[rgb]{0.81,0.36,0.00}{\textbf{#1}}}
\newcommand{\OtherTok}[1]{\textcolor[rgb]{0.56,0.35,0.01}{#1}}
\newcommand{\PreprocessorTok}[1]{\textcolor[rgb]{0.56,0.35,0.01}{\textit{#1}}}
\newcommand{\RegionMarkerTok}[1]{#1}
\newcommand{\SpecialCharTok}[1]{\textcolor[rgb]{0.81,0.36,0.00}{\textbf{#1}}}
\newcommand{\SpecialStringTok}[1]{\textcolor[rgb]{0.31,0.60,0.02}{#1}}
\newcommand{\StringTok}[1]{\textcolor[rgb]{0.31,0.60,0.02}{#1}}
\newcommand{\VariableTok}[1]{\textcolor[rgb]{0.00,0.00,0.00}{#1}}
\newcommand{\VerbatimStringTok}[1]{\textcolor[rgb]{0.31,0.60,0.02}{#1}}
\newcommand{\WarningTok}[1]{\textcolor[rgb]{0.56,0.35,0.01}{\textbf{\textit{#1}}}}
\usepackage{graphicx}
\makeatletter
\newsavebox\pandoc@box
\newcommand*\pandocbounded[1]{% scales image to fit in text height/width
  \sbox\pandoc@box{#1}%
  \Gscale@div\@tempa{\textheight}{\dimexpr\ht\pandoc@box+\dp\pandoc@box\relax}%
  \Gscale@div\@tempb{\linewidth}{\wd\pandoc@box}%
  \ifdim\@tempb\p@<\@tempa\p@\let\@tempa\@tempb\fi% select the smaller of both
  \ifdim\@tempa\p@<\p@\scalebox{\@tempa}{\usebox\pandoc@box}%
  \else\usebox{\pandoc@box}%
  \fi%
}
% Set default figure placement to htbp
\def\fps@figure{htbp}
\makeatother
\setlength{\emergencystretch}{3em} % prevent overfull lines
\providecommand{\tightlist}{%
  \setlength{\itemsep}{0pt}\setlength{\parskip}{0pt}}
\usepackage{amsmath}
\usepackage{amssymb}
\usepackage{booktabs}
\usepackage{graphicx}
\usepackage{float}
\usepackage{setspace}
\onehalfspacing
\usepackage{bookmark}
\IfFileExists{xurl.sty}{\usepackage{xurl}}{} % add URL line breaks if available
\urlstyle{same}
\hypersetup{
  pdftitle={Taller 2 -- Curva CO en planes de muestreo},
  pdfauthor={Diego Fernando Muñoz Portela 2415620, Valery Rivera Lopez 2417038},
  hidelinks,
  pdfcreator={LaTeX via pandoc}}

\title{Taller 2 -- Curva CO en planes de muestreo}
\author{Diego Fernando Muñoz Portela 2415620, Valery Rivera Lopez
2417038}
\date{02 de marzo de 2026}

\begin{document}
\maketitle

{
\setcounter{tocdepth}{3}
\tableofcontents
}
\section{Punto 1}\label{punto-1}

\subsection{Enunciado}\label{enunciado}

Suponga que usted tiene a su cargo la responsabilidad de recepción de
materia prima en una prestigiosa empresa de la región. En un acuerdo al
que se ha llegado con uno de sus proveedores, quien envía la materia
prima en lotes de 150 unidades, se estableció como un lote de muy buena
calidad aquel que presente un porcentaje de unidades defectuosas cercano
al 5\%.

Por otro lado el departamento de producción no está dispuesto a tolerar
ningún lote con más del 10\% de unidades defectuosas.

Como un lote se considera bueno si tiene menos del 5\%, se tomará un
tamaño de muestra \(n\) y se aceptará el lote si tiene menos del 5\% de
unidades defectuosas, por ejemplo:

\begin{itemize}
\tightlist
\item
  \((n=100, c=5)\)\\
\item
  \((n=80, c=4)\)\\
\item
  \((n=40, c=2)\)
\end{itemize}

\#\#Preguntas

\begin{enumerate}
\def\labelenumi{\alph{enumi}.}
\item
  Evalué en cada caso los riesgos del productor y consumidor.
\item
  ¿Se puede pensar entonces que un plan de muestreo con la condición de
  proporcionalidad en AQL es equivalente, sin importar el tamaño de
  muestra?.
\end{enumerate}

\begin{Shaded}
\begin{Highlighting}[]
\FunctionTok{library}\NormalTok{(tidyverse)}
\end{Highlighting}
\end{Shaded}

\begin{Shaded}
\begin{Highlighting}[]
\FunctionTok{library}\NormalTok{(ggthemes)}
\CommentTok{\# Definir los argumentos de LOS planes del muestreo }
\CommentTok{\#Fijos para los 3 planes}
\NormalTok{N}\OtherTok{=}\DecValTok{150}
\NormalTok{P }\OtherTok{\textless{}{-}} \FunctionTok{seq}\NormalTok{(}\FloatTok{0.000}\NormalTok{ , }\FloatTok{0.140}\NormalTok{, }\FloatTok{0.001}\NormalTok{)}

\CommentTok{\#PLAN 1}

\NormalTok{N }\OtherTok{\textless{}{-}} \DecValTok{150}
\NormalTok{n1 }\OtherTok{\textless{}{-}} \DecValTok{100}
\NormalTok{c1 }\OtherTok{\textless{}{-}} \DecValTok{5}

\CommentTok{\# Funcion para seleccionar la distribucion en cada uno de los planes }

\NormalTok{distribucion }\OtherTok{\textless{}{-}} \ControlFlowTok{function}\NormalTok{(n, N)\{}
\NormalTok{  r }\OtherTok{\textless{}{-}}\NormalTok{ n}\SpecialCharTok{/}\NormalTok{N}
  
  \ControlFlowTok{if}\NormalTok{(r }\SpecialCharTok{\textgreater{}} \FloatTok{0.1}\NormalTok{)\{}
    \FunctionTok{return}\NormalTok{(}\StringTok{"Hipergeométrica"}\NormalTok{)}
\NormalTok{  \}}
  
  \ControlFlowTok{if}\NormalTok{(r }\SpecialCharTok{\textless{}=} \FloatTok{0.1} \SpecialCharTok{\&\&}\NormalTok{ n}\SpecialCharTok{*}\NormalTok{r }\SpecialCharTok{\textgreater{}} \DecValTok{1}\NormalTok{)\{}
    \FunctionTok{return}\NormalTok{(}\StringTok{"Poisson"}\NormalTok{)}
\NormalTok{  \}}
  
  \ControlFlowTok{if}\NormalTok{(r }\SpecialCharTok{\textless{}=} \FloatTok{0.1}\NormalTok{)\{}
    \FunctionTok{return}\NormalTok{(}\StringTok{"Binomial"}\NormalTok{)}
\NormalTok{  \}}
\NormalTok{\}}

\CommentTok{\# Prueba distribución}
\FunctionTok{distribucion}\NormalTok{(n1, N)}
\end{Highlighting}
\end{Shaded}

\begin{verbatim}
## [1] "Hipergeométrica"
\end{verbatim}

\begin{Shaded}
\begin{Highlighting}[]
\CommentTok{\# Vector p}
\NormalTok{P }\OtherTok{\textless{}{-}} \FunctionTok{seq}\NormalTok{(}\FloatTok{0.000}\NormalTok{ , }\FloatTok{0.140}\NormalTok{, }\FloatTok{0.001}\NormalTok{)}

\CommentTok{\# Número de defectuosos en el lote}
\NormalTok{m }\OtherTok{\textless{}{-}} \FunctionTok{round}\NormalTok{(P }\SpecialCharTok{*}\NormalTok{ N)}

\CommentTok{\# Probabilidad de aceptación}
\NormalTok{Pa1 }\OtherTok{\textless{}{-}} \FunctionTok{phyper}\NormalTok{(c1, m, N}\SpecialCharTok{{-}}\NormalTok{m, n1)}

\CommentTok{\# Riesgos}
\NormalTok{AQL }\OtherTok{\textless{}{-}} \FloatTok{0.05}
\NormalTok{LTPD }\OtherTok{\textless{}{-}} \FloatTok{0.10}
\NormalTok{mAQL }\OtherTok{\textless{}{-}} \FunctionTok{round}\NormalTok{(AQL}\SpecialCharTok{*}\NormalTok{N)}
\NormalTok{mLTPD }\OtherTok{\textless{}{-}} \FunctionTok{round}\NormalTok{(LTPD}\SpecialCharTok{*}\NormalTok{N)}

\CommentTok{\# Riesgo productor}
\NormalTok{alpha1 }\OtherTok{\textless{}{-}} \DecValTok{1} \SpecialCharTok{{-}} \FunctionTok{phyper}\NormalTok{(c1, mAQL, N}\SpecialCharTok{{-}}\NormalTok{mAQL, n1)}
\NormalTok{alpha1}
\end{Highlighting}
\end{Shaded}

\begin{verbatim}
## [1] 0.4654046
\end{verbatim}

\begin{Shaded}
\begin{Highlighting}[]
\CommentTok{\# Riesgo consumidor}
\NormalTok{beta1 }\OtherTok{\textless{}{-}} \FunctionTok{phyper}\NormalTok{(c1, mLTPD, N}\SpecialCharTok{{-}}\NormalTok{mLTPD, n1)}
\NormalTok{beta1}
\end{Highlighting}
\end{Shaded}

\begin{verbatim}
## [1] 0.005796643
\end{verbatim}

\begin{Shaded}
\begin{Highlighting}[]
\CommentTok{\#CO1}
\FunctionTok{plot}\NormalTok{(P, Pa1, }\AttributeTok{type=}\StringTok{"l"}\NormalTok{, }\AttributeTok{lwd=}\DecValTok{3}\NormalTok{, }\AttributeTok{col=}\StringTok{"blue"}\NormalTok{,}
     \AttributeTok{xlab=}\StringTok{"Proporción defectuosa (p)"}\NormalTok{,}
     \AttributeTok{ylab=}\StringTok{"Probabilidad de aceptación (Pa)"}\NormalTok{,}
     \AttributeTok{main=}\StringTok{"Curva CO {-} Plan 1 (n=100, c=5)"}\NormalTok{,}
     \AttributeTok{ylim=}\FunctionTok{c}\NormalTok{(}\DecValTok{0}\NormalTok{,}\DecValTok{1}\NormalTok{))}
\FunctionTok{grid}\NormalTok{()}
\end{Highlighting}
\end{Shaded}

\begin{center}\includegraphics{RMK_FINAL_files/figure-latex/plan-1} \end{center}

\section{Punto 2}\label{punto-2}

\subsection{Enunciado}\label{enunciado-1}

Construya una Curva característica de operación para un plan de muestreo
simple en el cual el tamaño del lote es de 100 unidades, tamaño de
muestra 50, y numero de aceptación 2. Utilice las tres distribuciones de
probabilidad.

\begin{enumerate}
\def\labelenumi{\alph{enumi}.}
\tightlist
\item
  ¿Se presentan discrepancias entre las curvas?
\end{enumerate}

\begin{Shaded}
\begin{Highlighting}[]
\CommentTok{\# Parámetros}
\NormalTok{N }\OtherTok{\textless{}{-}} \DecValTok{100}
\NormalTok{n }\OtherTok{\textless{}{-}} \DecValTok{50}
\NormalTok{c }\OtherTok{\textless{}{-}} \DecValTok{2}

\CommentTok{\# Vector de proporciones defectuosas}
\NormalTok{P }\OtherTok{\textless{}{-}} \FunctionTok{seq}\NormalTok{(}\FloatTok{0.001}\NormalTok{, }\FloatTok{0.20}\NormalTok{, }\FloatTok{0.001}\NormalTok{)}

\CommentTok{\# 1. Binomial}
\NormalTok{Pa\_bin }\OtherTok{\textless{}{-}} \FunctionTok{pbinom}\NormalTok{(c, n, P)}

\CommentTok{\# 2. Poisson}
\NormalTok{Pa\_pois }\OtherTok{\textless{}{-}} \FunctionTok{ppois}\NormalTok{(c, n}\SpecialCharTok{*}\NormalTok{P)}

\CommentTok{\# 3. Hipergeométrica}
\NormalTok{Pa\_hiper }\OtherTok{\textless{}{-}} \FunctionTok{sapply}\NormalTok{(P, }\ControlFlowTok{function}\NormalTok{(p)\{}
\NormalTok{  D }\OtherTok{\textless{}{-}} \FunctionTok{round}\NormalTok{(N}\SpecialCharTok{*}\NormalTok{p)}
  \FunctionTok{phyper}\NormalTok{(c, D, N}\SpecialCharTok{{-}}\NormalTok{D, n)}
\NormalTok{\})}

\CommentTok{\# {-}{-}{-} GRÁFICA MEJORADA {-}{-}{-}}

\FunctionTok{par}\NormalTok{(}\AttributeTok{bg=}\StringTok{"white"}\NormalTok{)  }\CommentTok{\# fondo blanco}

\FunctionTok{plot}\NormalTok{(P, Pa\_hiper, }\AttributeTok{type=}\StringTok{"l"}\NormalTok{, }\AttributeTok{lwd=}\DecValTok{3}\NormalTok{, }\AttributeTok{col=}\StringTok{"darkgreen"}\NormalTok{,}
     \AttributeTok{ylab=}\StringTok{"Probabilidad de aceptación (Pa)"}\NormalTok{,}
     \AttributeTok{xlab=}\StringTok{"Proporción defectuosa (p)"}\NormalTok{,}
     \AttributeTok{main=}\StringTok{"Curva Característica de Operación}\SpecialCharTok{\textbackslash{}n}\StringTok{N=100, n=50, c=2"}\NormalTok{,}
     \AttributeTok{ylim=}\FunctionTok{c}\NormalTok{(}\DecValTok{0}\NormalTok{,}\DecValTok{1}\NormalTok{))}

\FunctionTok{grid}\NormalTok{(}\AttributeTok{col=}\StringTok{"gray85"}\NormalTok{)}

\FunctionTok{lines}\NormalTok{(P, Pa\_bin, }\AttributeTok{lwd=}\DecValTok{3}\NormalTok{, }\AttributeTok{lty=}\DecValTok{2}\NormalTok{, }\AttributeTok{col=}\StringTok{"blue"}\NormalTok{)}
\FunctionTok{lines}\NormalTok{(P, Pa\_pois, }\AttributeTok{lwd=}\DecValTok{3}\NormalTok{, }\AttributeTok{lty=}\DecValTok{3}\NormalTok{, }\AttributeTok{col=}\StringTok{"red"}\NormalTok{)}

\FunctionTok{legend}\NormalTok{(}\StringTok{"topright"}\NormalTok{,}
       \AttributeTok{legend=}\FunctionTok{c}\NormalTok{(}\StringTok{"Hipergeométrica"}\NormalTok{,}\StringTok{"Binomial"}\NormalTok{,}\StringTok{"Poisson"}\NormalTok{),}
       \AttributeTok{col=}\FunctionTok{c}\NormalTok{(}\StringTok{"darkgreen"}\NormalTok{,}\StringTok{"blue"}\NormalTok{,}\StringTok{"red"}\NormalTok{),}
       \AttributeTok{lwd=}\DecValTok{3}\NormalTok{,}
       \AttributeTok{lty=}\FunctionTok{c}\NormalTok{(}\DecValTok{1}\NormalTok{,}\DecValTok{2}\NormalTok{,}\DecValTok{3}\NormalTok{),}
       \AttributeTok{bty=}\StringTok{"n"}\NormalTok{)}
\end{Highlighting}
\end{Shaded}

\begin{center}\includegraphics{RMK_FINAL_files/figure-latex/plan23-1} \end{center}

La distribución hipergeométrica (verde) presenta un comportamiento
escalonado debido a su naturaleza discreta y al tamaño finito del lote.
Dado que 𝑁 = 100 N=100, los cambios en el número de defectuosos 𝐷 = 𝑁 𝑝
D=Np ocurren en valores enteros, generando variaciones abruptas en la
probabilidad de aceptación.

Además, esta distribución describe con mayor precisión el riesgo real
del plan de muestreo, ya que considera el agotamiento de la población
durante el muestreo (sin reemplazo).

Por el contrario, las aproximaciones binomial y Poisson asumen
poblaciones infinitas o muestreo con reemplazo, lo que genera ligeras
sobreestimaciones en la probabilidad de aceptación.

\begin{enumerate}
\def\labelenumi{\alph{enumi}.}
\setcounter{enumi}{1}
\tightlist
\item
  ¿Que sucede si el tamaño del lote es de 2000 unidades?
\end{enumerate}

\begin{Shaded}
\begin{Highlighting}[]
\CommentTok{\# Parámetros}
\NormalTok{N }\OtherTok{\textless{}{-}} \DecValTok{2000}
\NormalTok{n }\OtherTok{\textless{}{-}} \DecValTok{50}
\NormalTok{c }\OtherTok{\textless{}{-}} \DecValTok{2}

\CommentTok{\# Vector de proporciones defectuosas}
\NormalTok{P }\OtherTok{\textless{}{-}} \FunctionTok{seq}\NormalTok{(}\FloatTok{0.001}\NormalTok{, }\FloatTok{0.20}\NormalTok{, }\FloatTok{0.001}\NormalTok{)}

\CommentTok{\# 1. Binomial}
\NormalTok{Pa\_bin }\OtherTok{\textless{}{-}} \FunctionTok{pbinom}\NormalTok{(c, n, P)}

\CommentTok{\# 2. Poisson}
\NormalTok{Pa\_pois }\OtherTok{\textless{}{-}} \FunctionTok{ppois}\NormalTok{(c, n}\SpecialCharTok{*}\NormalTok{P)}

\CommentTok{\# 3. Hipergeométrica}
\NormalTok{Pa\_hiper }\OtherTok{\textless{}{-}} \FunctionTok{sapply}\NormalTok{(P, }\ControlFlowTok{function}\NormalTok{(p)\{}
\NormalTok{  D }\OtherTok{\textless{}{-}} \FunctionTok{round}\NormalTok{(N}\SpecialCharTok{*}\NormalTok{p)}
  \FunctionTok{phyper}\NormalTok{(c, D, N}\SpecialCharTok{{-}}\NormalTok{D, n)}
\NormalTok{\})}

\CommentTok{\# {-}{-}{-} GRÁFICA MEJORADA {-}{-}{-}}

\FunctionTok{par}\NormalTok{(}\AttributeTok{bg=}\StringTok{"white"}\NormalTok{)  }\CommentTok{\# fondo blanco}

\FunctionTok{plot}\NormalTok{(P, Pa\_hiper, }\AttributeTok{type=}\StringTok{"l"}\NormalTok{, }\AttributeTok{lwd=}\DecValTok{3}\NormalTok{, }\AttributeTok{col=}\StringTok{"darkgreen"}\NormalTok{,}
     \AttributeTok{ylab=}\StringTok{"Probabilidad de aceptación (Pa)"}\NormalTok{,}
     \AttributeTok{xlab=}\StringTok{"Proporción defectuosa (p)"}\NormalTok{,}
     \AttributeTok{main=}\StringTok{"Curva Característica de Operación}\SpecialCharTok{\textbackslash{}n}\StringTok{N=2000, n=50, c=2"}\NormalTok{,}
     \AttributeTok{ylim=}\FunctionTok{c}\NormalTok{(}\DecValTok{0}\NormalTok{,}\DecValTok{1}\NormalTok{))}

\FunctionTok{grid}\NormalTok{(}\AttributeTok{col=}\StringTok{"gray85"}\NormalTok{)}

\FunctionTok{lines}\NormalTok{(P, Pa\_bin, }\AttributeTok{lwd=}\DecValTok{3}\NormalTok{, }\AttributeTok{lty=}\DecValTok{2}\NormalTok{, }\AttributeTok{col=}\StringTok{"blue"}\NormalTok{)}
\FunctionTok{lines}\NormalTok{(P, Pa\_pois, }\AttributeTok{lwd=}\DecValTok{3}\NormalTok{, }\AttributeTok{lty=}\DecValTok{3}\NormalTok{, }\AttributeTok{col=}\StringTok{"red"}\NormalTok{)}

\FunctionTok{legend}\NormalTok{(}\StringTok{"topright"}\NormalTok{,}
       \AttributeTok{legend=}\FunctionTok{c}\NormalTok{(}\StringTok{"Hipergeométrica"}\NormalTok{,}\StringTok{"Binomial"}\NormalTok{,}\StringTok{"Poisson"}\NormalTok{),}
       \AttributeTok{col=}\FunctionTok{c}\NormalTok{(}\StringTok{"darkgreen"}\NormalTok{,}\StringTok{"blue"}\NormalTok{,}\StringTok{"red"}\NormalTok{),}
       \AttributeTok{lwd=}\DecValTok{3}\NormalTok{,}
       \AttributeTok{lty=}\FunctionTok{c}\NormalTok{(}\DecValTok{1}\NormalTok{,}\DecValTok{2}\NormalTok{,}\DecValTok{3}\NormalTok{),}
       \AttributeTok{bty=}\StringTok{"n"}\NormalTok{)}
\end{Highlighting}
\end{Shaded}

\begin{center}\includegraphics{RMK_FINAL_files/figure-latex/plan24-1} \end{center}

\begin{center}\rule{0.5\linewidth}{0.5pt}\end{center}

\end{document}
